Диссертационная работа посвящена разработке методов и алгоритмов
анализа консистентности памяти многопроцессорных систем с общей
памятью при исполнении многопоточных программ. В работе
рассматривается задача верификации аппаратной реализации
подсистемы памяти с учётом модели памяти в условиях
недетерминированного порядка операций.
В основе исследования лежит представление многопоточных программ в
виде частично упорядоченных множеств операций с памятью. Особое
внимание уделено разработке эффективных алгоритмов построения
линейных порядков таких множеств с учётом классов эквивалентности
перестановок, а также созданию структуры данных для полного описания
всех возможных результатов исполнения программы, объединяющей
эквивалентные промежуточные состояния.
Предложенный подход не накладывает ограничений на операции с памятью,
что обеспечивает его применимость в широком спектре тестовых сценариев.

{\actuality} связана с ростом числа ядер в многопроцессорных
системах, что приводит к существенному усложнению архитектуры
подсистемы памяти и, как следствие, к значительным трудностям при её
верификации.
Рост тактовой частоты прекратился, хотя ранее он служил основным
фактором повышения производительности при разработке процессоров
общего назначения~\cite{cpu,trend}.
По пути повышения быстродействия вычислительное ядро приблизилось к
своим физическим и технологическим пределам~\cite{physLimits}.
Индустрия обратилась к новым стратегиям для дальнейшего роста
производительности, и ключевым направлением развития стала
многоядерность~\cite{multithreaded}.
Данное решение позволило организовать параллельное выполнение задач
на процессоре,
увеличив общую эффективность вычислительной системы.
Вследствие этого сложность таких систем значительно увеличилась.
В них применяются многоуровневые иерархические архитектуры,
параллелизм исполнения задач требует сложных механизмов
синхронизации~\cite{sync},
усложняя взаимодействие между различными компонентами~\cite{hieracy}.
Рост сложности современных вычислительных систем опережает развитие
языковых средств их описания.
Несмотря на широкое распространение в промышленности языков описания
аппаратуры (Hardware Description Language,
HDL)~\cite{mermet2012fundamentals}, таких как VHDL~\cite{vhdl},
Verilog~\cite{verilog}
и его расширенная версия
SystemVerilog~\cite{systemvetilog}, процессы верификации, анализа и
отладки моделей, создаваемых на них, остаются крайне трудоемкими.
Отчасти данную проблему решают высокоуровневые языки описания
архитектуры (High-Level Synthesis, HLS)~\cite{hls}, которые предоставляют
разработчикам современные парадигмы проектирования, более высокий
уровень абстракции, а также встроенные возможности для анализа и валидации кода.
Однако даже с их использованием обеспечение полной верификации
создаваемой аппаратуры на соответствие спецификации представляется
сложной задачей ввиду общей сложности
современных вычислительных систем.

Одним из наиболее сложных аспектов такой
верификации является проверка консистентности памяти -- корректности
результатов операций чтения и записи при параллельном исполнении программ.

В многопроцессорных системах~\cite{smp} с общей памятью модель
памяти~\cite{memory-model} -- часть архитектуры процессора -- определяет
порядок, в котором результаты исполнения операций из разных потоков
становятся видимыми друг для друга.

Нарушение спецификации может приводить к трудно
воспроизводимым сбоям и ошибкам, проявляющимся лишь в
определённых сценариях выполнения.
Таким образом, разработка эффективных методов анализа консистентности
памяти с учётом моделей памяти представляет собой актуальную задачу,
рассматриваемую в данном исследовании.

{\progress}.
Теоретическую основу для описания поведения параллельных систем даёт
теория трасс~\cite{mazurkiewicz1986trace,mazurkiewicz1995introduction},
в которой задаётся алфавит действий системы и отношение
зависимости между ними.
Пары действий, являющиеся независимыми, могут
происходить в любом порядке без изменения результата.
Строки, отличающиеся лишь перестановкой соседних независимых действий,
относятся к одному классу эквивалентности -- трассе, что позволяет
сокращать пространство поиска при анализе конкурентного исполнения.
На практике для верификации моделей памяти используются алгоритмы
частичного и динамического сокращения порядка
(POR/DPOR)~\cite{flanagan2005dynamic,chockler2018computer},
которые исследуют ровно по одному представителю от каждой трассы и
тем самым позволяют избегать перебора всех допустимых чередований.
Однако, поскольку подобные алгоритмы оперируют лишь отношением
зависимости, среди порождаемых кандидатов могут присутствовать
исполнения, исключаемые лишь на последующих этапах проверки аксиом
выбранной модели памяти~\cite{clarke1997model}.
Этот подход применяется в инструментах, порождающих возможные трассы
исполнения, в частности, в литмус-тестировании~\cite{litmus2},
которое служит отправной точкой для исследования эффектов различных
моделей памяти. Его цель -- выявление отличительных особенностей
моделей и проверка того, соответствует ли поведение системы заданной
модели памяти. Для этого используется значительное число уникальных
тестовых сценариев, успешное прохождение которых подтверждает базовую
работоспособность многопроцессорных систем

Другое направление связано с проверкой консистентности памяти в
условиях большого числа конкурирующих
обращений~\cite{manovit, Axe}. В отличие от литмус-тестирования,
где каждый тест содержит не более десяти операций, данные подходы
оперируют тысячами и десятками тысяч обращений к памяти, что
позволяет выявлять ошибки, проявляющиеся только при интенсивном
взаимодействии ядер через общую память. Однако увеличение числа
операций влечёт за собой экспоненциальный рост количества возможных
трасс исполнения, что требует применения эвристик и ограничений на
структуру тестов.

Описанные выше подходы можно отнести к методам верификации подсистемы
памяти на готовом кристалле
(post-silicon)~\cite{post-silicon-book}, при которых тестирование не
привязано к деталям реализации многопроцессорной системы.
Наряду с ними до изготовления кристалла (pre-silicon) широко
применяются другие методы верификации
~\cite{andrade2019chaining, andrade2020directed, breker1}.
Их отличительной особенностью является направленное тестирование
отдельных состояний подсистемы памяти. Данные методы используются на
ранних стадиях разработки, однако зачастую не обладают
масштабируемостью, а используемые в них инструменты являются проприетарными.
Отдельного внимания заслуживают механизмы верификации, интегрированные
непосредственно в подсистему памяти~\cite{post-silicon}. В отличие от
рассмотренных выше подходов, данный класс решений предполагает
встраивание логики трассировки обращений
непосредственно в аппаратную часть системы. Основным недостатком
такого подхода является необходимость аппаратной реализации
дополнительной логики, что требует ресурсов на кристалле,
увеличивает стоимость разработки и делает подход неприменимым для
верификации систем, не имеющих такой поддержки.

Настоящая работа сосредоточена на post-silicon подходах к проверке
подсистемы памяти, не зависящих от деталей реализации системы.
Классические подходы к анализу конкурентного исполнения
используют единое отношение зависимости для определения как
допустимых перестановок, так и их эквивалентности.
В результате среди порождаемых кандидатов могут присутствовать
исполнения, нарушающие аксиомы выбранной модели памяти.
Либо используются подходы, с существенными ограничениями на структуру теста.
В данной работе предлагается разделить ограничения допустимости
порядка и зависимость операций.
Частично упорядоченное множество $P$ (partially ordered set)~\cite{poset}
операций с памятью задаёт ограничения, накладываемые моделью памяти,
тогда как $D$ представляет собой отношение зависимости,
определяющее значимые варианты взаимного порядка операций.

{\aim} данной работы является создание методов и разработка
алгоритмов анализа консистентности памяти при исполнении
многопоточных программ в многопроцессорных системах с учётом моделей памяти.

Для~достижения поставленной цели необходимо было решить следующие {\tasks}:
\begin{enumerate}[beginpenalty=10000] %
    % https://tex.stackexchange.com/a/476052/104425
  \item сведение многопоточной программы к частично упорядоченному
    множеству операций с памятью;
  \item разработка алгоритмов построения линейных порядков
    частично упорядоченного множества, учитывающих свойства
    эквивалентных перестановок;
  \item разработка структуры данных для полного описания результатов
    исполнения многопоточной программы, объединяющей эквивалентные
    промежуточные состояния (графа состояний программы);
  \item разработка алгоритмов верификации соответствия аппаратуры
    заявленной модели памяти на основе анализа результатов исполнения
    многопоточной программы.
\end{enumerate}

Тема и содержание диссертационной работы соответствуют паспорту
научной специальности 2.3.5. Математическое и программное обеспечение
вычислительных систем, комплексов и компьютерных сетей, в
частности пунктам:
\begin{itemize}
  \item п. 1 -- Модели, методы и алгоритмы проектирования, анализа,
    трансформации, верификации и тестирования программ и программных систем.
  \item п. 8 -- Модели и методы создания программ и программных
    систем для параллельной и распределенной обработки данных, языки
    и инструментальные средства параллельного программирования.
\end{itemize}

% Объект исследования представляет собой знание, порождающее проблемную ситуацию
{\object} является подсистема памяти многопроцессорных систем с общей памятью.

% Предметом исследования являются технические решения, методы,
% алгоритмы, технологии для
{\subject} являются многопоточные программы и операции с памятью в
них, математические модели, алгоритмы и структуры
данных для анализа консистентности памяти многопроцессорных систем, а
именно: методы представления многопоточных программ в виде частично
упорядоченных множеств операций с памятью, алгоритмы построения
линейных порядков с учётом классов эквивалентности, а также граф
состояний как средство представления результатов исполнения.

{\novelty}. В данной работе впервые:
\begin{enumerate}[beginpenalty=10000] %
    % How to prevent a page break before an itemize list?
    % https://tex.stackexchange.com/a/476052/104425
  \item предложена математическая модель анализа многопоточной
    программы, в которой ограничения, задаваемые моделью памяти,
    представляются частичным порядком $P$ над операциями с памятью, а
    значимые варианты взаимного порядка операций --
    отношением зависимости $D$;
  \item разработаны алгоритмы построения одного представителя
    каждого класса эквивалентности линейных порядков заданного
    частичного порядка $P$;
  \item предложена структура данных -- граф состояний, объединяющий
    эквивалентные промежуточные состояния при анализе различных
    линейных порядков и тем самым обеспечивающий
    представление множества возможных результатов без явного хранения
    всех линейных порядков;
  \item предложены алгоритмы верификации результата исполнения
    программы, как на основе линейных порядков частично
    упорядоченного множества, так и на
    основе графа состояний.
\end{enumerate}

{\influence}.
Теоретическая значимость работы состоит в развитии методов анализа
многопоточных программ на основе частичных порядков и классов
эквивалентности перестановок. Раздельное представление ограничений,
накладываемых моделью памяти на допустимые порядки операций, и
отношения зависимости операций позволяет рассматривать пространство
допустимых исполнений как множество классов эквивалентности линейных
порядков частичного порядка.
Для предложенной математической модели доказаны корректность и
полнота алгоритмов построения представителей классов эквивалентности
без повторного порождения представителей одного класса. Обоснована
возможность полного представления множества допустимых результатов
исполнения графом состояний, объединяющим эквивалентные промежуточные
состояния. Доказана эквивалентность методов проверки допустимости
результатов исполнения на основе представителей классов линейных
порядков и графа состояний.
Практическая значимость научных результатов заключается в разработке
инструментария верификации подсистемы памяти, снимающего ограничения
существующих подходов в отношении типов операций и пересечений
адресов, а также применимого к многопоточным программам с большим числом
конкурирующих обращений к памяти. Предложенные методы могут быть
направлены на выявление критических ошибок на
ранних этапах разработки многопроцессорных систем, что способствует повышению
их отказоустойчивости и стабильности.

Экспериментальная апробация разработанных алгоритмов проводилась на
тестовых сценариях, моделирующих как литмус-тесты, так и нагрузки с
большим числом конкурирующих обращений к памяти.
Результаты работы были использованы в \fixme{where?} в рамках
приёмочного тестирования, что подтверждает их практическую применимость
и востребованность.

{\methods}.
Для решения поставленных задач в работе используется
методы дискретной математики и теории алгоритмов, включая
теорию частично упорядоченных множеств, комбинаторный анализ, теорию
графов.
Также в практической части работы применялись методы планирования эксперимента,
измерения, анализа и сравнения результатов.

{\defpositions}:
\begin{enumerate}[beginpenalty=10000]
  \item раздельное представление ограничений на порядок операций с
    памятью в виде частичного порядка $P$ и зависимостей между
    операциями в виде отношения $D$ позволяет описывать множество
    допустимых исполнений многопоточной программы через классы
    эквивалентности линейных расширений $P$ и исключать порядки,
    нарушающие аксиомы модели памяти, до их явного построения;
  \item разработанные алгоритмы перечисления линейных порядков
    частичного порядка $P$ с учётом отношения зависимости $D$
    обеспечивают построение ровно одного представителя каждого класса
    эквивалентности, гарантируя полноту перечисления и отсутствие повторений;
  \item предложенный граф состояний обеспечивает полное представление
    множества возможных результатов исполнения многопоточной
    программы без явного хранения всех линейных порядков, за счёт
    объединения эквивалентных промежуточных состояний;
  \item доказана применимость и эквивалентность алгоритмов
    верификации результатов исполнения многопоточных программ на
    основе линейных порядков и графа состояний.
\end{enumerate}

{\reliability}
научных результатов, представленных в диссертационной работе, обеспечивается
строгостью построений. Все ключевые утверждения, включая
условия корректности исполнения многопоточных программ и
доказательства эквивалентности подходов к верификации,
обоснованы в тексте работы.
Разработанные алгоритмы подтверждены обоснованными оценками
асимптотической сложности, которые согласуются с известными
теоретическими результатами в комбинаторике и теории сложности вычислений.
Кроме того, достоверность подтверждается экспериментальным
тестированием на эталонных сценариях, которые моделируют как
классические литмус-тесты, так и нагрузки с большим
количеством конкурирующих обращений к памяти, что демонстрирует
практическую применимость и воспроизводимость полученных результатов.
Основные положения и результаты диссертационной работы были
апробированы в ходе докладов на научных конференциях.

{\probation}.
Основные результаты работы докладывались~на:
\begin{enumerate}
  \item 68-й Всероссийской научной конференции МФТИ в честь 130-летия
    со дня рождения Н.Н. Семенова:
    (30 марта - 04 апреля 2026г.)~\cite{mipt68}\fixme{fix ref and
    disp when it is published}
\end{enumerate}

% {\contribution} Автор принимал активное участие \ldots

\ifnumequal{\value{bibliosel}}{0}
{%%% Встроенная реализация с загрузкой файла через движок bibtex8.
  % (При желании, внутри можно использовать обычные ссылки, наподобие
  % `\cite{vakbib1,vakbib2}`).
  {\publications}.
  Основные результаты по теме диссертации изложены
  в~XX~печатных изданиях,
  X из которых изданы в журналах, рекомендованных ВАК,
  X "--- в тезисах докладов.
}%
{%%% Реализация пакетом biblatex через движок biber
  \begin{refsection}[bl-author, bl-registered]
    % Это refsection=1.
    % Процитированные здесь работы:
    %  * подсчитываются, для автоматического составления фразы
    % "Основные результаты ..."
    %  * попадают в авторскую библиографию, при usefootcite==0 и
    % стиле `\insertbiblioauthor` или `\insertbiblioauthorgrouped`
    %  * нумеруются там в зависимости от порядка команд
    % `\printbibliography` в этом разделе.
    %  * при использовании `\insertbiblioauthorgrouped`, порядок
    % команд `\printbibliography` в нём должен быть тем же (см.
    % biblio/biblatex.tex)
    %
    % Невидимый библиографический список для подсчёта количества публикаций:
    \phantom{\printbibliography[heading=nobibheading, section=1,
      env=countauthorvak,          keyword=biblioauthorvak]%
      \printbibliography[heading=nobibheading, section=1,
      env=countauthorwos,          keyword=biblioauthorwos]%
      \printbibliography[heading=nobibheading, section=1,
      env=countauthorscopus,       keyword=biblioauthorscopus]%
      \printbibliography[heading=nobibheading, section=1,
      env=countauthorconf,         keyword=biblioauthorconf]%
      \printbibliography[heading=nobibheading, section=1,
      env=countauthorother,        keyword=biblioauthorother]%
      \printbibliography[heading=nobibheading, section=1,
      env=countregistered,         keyword=biblioregistered]%
      \printbibliography[heading=nobibheading, section=1,
      env=countauthorpatent,       keyword=biblioauthorpatent]%
      \printbibliography[heading=nobibheading, section=1,
      env=countauthorprogram,      keyword=biblioauthorprogram]%
      \printbibliography[heading=nobibheading, section=1,
      env=countauthor,             keyword=biblioauthor]%
      \printbibliography[heading=nobibheading, section=1,
      env=countauthorvakscopuswos, filter=vakscopuswos]%
      \printbibliography[heading=nobibheading, section=1,
    env=countauthorscopuswos,    filter=scopuswos]}%
    %
    \nocite{*}%
    %
    {\publications} Основные результаты по теме диссертации изложены
    в~\arabic{citeauthor}~печатных изданиях,
    \arabic{citeauthorvak} из которых изданы в журналах,
    рекомендованных ВАК%
    \ifnum \value{citeauthorscopuswos}>0%
    , \arabic{citeauthorscopuswos} "--- в~периодических научных
    журналах, индексируемых Web of~Science и Scopus%
    \fi%
    \ifnum \value{citeauthorconf}>0%
    , \arabic{citeauthorconf} "--- в~тезисах докладов.
    \else%
    .
    \fi%
    \ifnum \value{citeregistered}=1%
    \ifnum \value{citeauthorpatent}=1%
    Зарегистрирован \arabic{citeauthorpatent} патент.
    \fi%
    \ifnum \value{citeauthorprogram}=1%
    Зарегистрирована \arabic{citeauthorprogram} программа для ЭВМ.
    \fi%
    \fi%
    \ifnum \value{citeregistered}>1%
    Зарегистрированы\ %
    \ifnum \value{citeauthorpatent}>0%
    \formbytotal{citeauthorpatent}{патент}{}{а}{}%
    \ifnum \value{citeauthorprogram}=0 . \else \ и~\fi%
    \fi%
    \ifnum \value{citeauthorprogram}>0%
    \formbytotal{citeauthorprogram}{программ}{а}{ы}{} для ЭВМ.
    \fi%
    \fi%
    % К публикациям, в которых излагаются основные научные результаты
    % диссертации на соискание учёной
    % степени, в рецензируемых изданиях приравниваются патенты на
    % изобретения, патенты (свидетельства) на
    % полезную модель, патенты на промышленный образец, патенты на
    % селекционные достижения, свидетельства
    % на программу для электронных вычислительных машин, базу данных,
    % топологию интегральных микросхем,
    % зарегистрированные в установленном порядке.(в ред.
    % Постановления Правительства РФ от 21.04.2016 N 335)
  \end{refsection}%
  \begin{refsection}[bl-author, bl-registered]
    % Это refsection=2.
    % Процитированные здесь работы:
    %  * попадают в авторскую библиографию, при usefootcite==0 и
    % стиле `\insertbiblioauthorimportant`.
    %  * ни на что не влияют в противном случае
    \nocite{sitito_1}
    \nocite{sitito_2}
    \nocite{mipt68}
    %\nocite{vakbib2}%vak
    %\nocite{patbib1}%patent
    %\nocite{progbib1}%program
    %\nocite{bib1}%other
    %\nocite{confbib1}%conf
  \end{refsection}%
  %
  % Всё, что вне этих двух refsection, это refsection=0,
  %  * для диссертации - это нормальные ссылки, попадающие в обычную
  % библиографию
  %  * для автореферата:
  %     * при usefootcite==0, ссылка корректно сработает только для
  % источника из `external.bib`. Для своих работ --- напечатает "[0]"
  % (и даже Warning не вылезет).
  %     * при usefootcite==1, ссылка сработает нормально. В авторской
  % библиографии будут только процитированные в refsection=0 работы.
}
